\section{Volume Hessians, finite elements, and computational complexity}
\label{sec:ipolybridge}

The variational description of a polyhedron has a second interpretation:
its infinitesimal deformation equation can be a discrete Dirichlet
problem. This is an exact connection to finite elements and electrical
networks, and it supplies the computational counterpart of the volume
and symplectic constructions developed above. The companion study
\emph{Duffin's resistor networks, the volume of ideal polyhedra, and
nearly-linear-time Laplacian solvers}~\cite{IpolyStudy} develops this
connection and its numerical examples. We incorporate the essential
identities here and connect them to boundary response, Gaussian
integration, flips, and the Casson--Rivin program.

\subsection{The cotangent form and the correct deformation coordinates}

Let $\mathcal T$ triangulate a planar disk, with interior vertices $I$
and boundary vertices $B$. For a piecewise affine function $v$ on a
geometric realization, integration over each triangle gives
\begin{equation}
 \int_{|\mathcal T|}|\nabla v|^2\,dA
 =\sum_{e=ij}w_e(v_i-v_j)^2=v^TLv,
 \qquad
 w_e=\frac12(\cot\alpha_e+\cot\beta_e).
 \label{eq:ipolyFEM}
\end{equation}
Here $\alpha_e,\beta_e$ are the angles opposite $e$, and the second
term is omitted on a boundary edge. To check the normalization, the
off-diagonal stiffness entry on a triangle is
$\int_T\nabla\phi_i\cdot\nabla\phi_j\,dA=-\tfrac12\cot\alpha_k$
for its barycentric functions. Row sums vanish, giving the displayed
edge form. This is the cotangent network of the finite element and
electrical interpretations~\cite{DuffinNetworks}.

The energy is nonnegative on every geometric mesh, even if some edge
weights are negative. The additional condition for nonnegative
\emph{interior} conductances is local Delaunayness:
\begin{equation}
 w_e=\frac{\sin(\alpha_e+\beta_e)}
              {2\sin\alpha_e\sin\beta_e}\ge0
 \quad\Longleftrightarrow\quad \alpha_e+\beta_e\le\pi.
 \label{eq:ipolyDelaunayweight}
\end{equation}
When the relevant weights are nonnegative, the grounded matrix $L_{II}$
is a positive definite graph Laplacian, and solving a Dirichlet problem
is solving Kirchhoff's equations. Boundary-to-boundary edges do not enter
this Dirichlet system; positivity of their weights is an additional
condition if the full matrix is used in a different boundary problem.

Now prescribe the sums $\Delta_e=\alpha_e+\beta_e\le\pi$ opposite
interior edges, and prescribe the single opposite angle on each
boundary edge. An angle structure assigns positive corner angles,
with sum $\pi$ in every triangle and these prescribed edge sums.
Assume this positive angle domain is nonempty. For planar realization
the data also satisfy
$\sum_{e\ni i}(\pi-\Delta_e)=2\pi$ at each interior vertex; the
left-hand side is its total angle in every such angle structure.
If the intended planar boundary is strictly convex, impose also
\[
 0<\kappa_b:=\sum_{t\ni b}\alpha_{t,b}
   =\sum_{e\ni b}(\pi-\Delta_e)-\pi<\pi
 \qquad(b\in B).
\]
Here $\Delta_e$ denotes the single prescribed opposite angle on a
boundary edge. These boundary inequalities are needed for the convex
polyhedral interpretation; the Hessian identities below already hold
on the abstract positive angle domain. Positive feasibility alone
does not impose convexity of the developed boundary.
Orient the triangles consistently. If a corner $c$ is opposite the
directed edge $i\to j$, set
\begin{equation}
 \alpha_c(u)=\alpha_c^0+u_i-u_j,\qquad u|_B=0.
 \label{eq:ipolyAngleCoordinates}
\end{equation}
These are leading--trailing coordinates, not logarithmic scale factors.
Triangle sums telescope; increments opposite a shared edge cancel.
For a triangulated disk they parametrize the entire affine space of
solutions: the map is injective, and the connected triangle--edge
incidence graph together with Euler's formula gives dimension $|I|$.

Let $\Lambda'(a)=-\log(2\sin a)$ be the Lobachevsky derivative and put
$\mathcal V(\alpha)=\sum_c\Lambda(\alpha_c)$. At a geometric
configuration this is the sum of the volumes of the ideal tetrahedra
obtained by adjoining the vertex at infinity to the planar triangles.
Rivin's variational principle uses this concave functional to impose
geometric closure~\cite{RivinEuclidean,FuterGueritaud}.

\begin{proposition}[The exact finite-element Hessian]
\label{prop:ipolyHessian}
On the positive angle domain in the coordinates
\eqref{eq:ipolyAngleCoordinates},
\begin{equation}
 \nabla_u\mathcal V=(\log H_i)_{i\in I},\qquad
 -D_u^2\mathcal V=2L_{II}(\alpha),
 \label{eq:ipolyHessian}
\end{equation}
where $H_i$ is the dilation accumulated by developing the incident
triangles once around $i$. In particular, at a geometric critical point
the negative volume Hessian is twice the finite element stiffness matrix.
\end{proposition}

\begin{proof}
In a positively oriented triangle $ijk$, differentiation with respect
to $u_i$ gives
$\log(\sin\alpha_{t,j}/\sin\alpha_{t,k})$.
Summing over the star of $i$ gives the logarithm of the product of
side ratios, which is the dilation holonomy by the sine law.
Since $\Lambda''(a)=-\cot a$, the two corners opposite an interior
edge contribute
$-(\cot\alpha_e+\cot\beta_e)(du_i-du_j)^2$ to $D^2\mathcal V$.
This is $-2$ times the corresponding summand of
\eqref{eq:ipolyFEM}. Boundary increments vanish.
\end{proof}

For a true interior edge with $0<\Delta_e<\pi$, the weights satisfy
\begin{equation}
 w_e=\frac{\sin\Delta_e}{2\sin\alpha_e\sin\beta_e}
       \ge\cot(\Delta_e/2)>0.
 \label{eq:ipolyWeightLower}
\end{equation}
Thus positivity holds throughout the angle domain, before the angles
close to form a geometric mesh. At $\Delta_e=\pi$ the weight is zero;
such a diagonal subdivides a flat face. The finite element interpretation
as an actual planar mesh applies at realization, whereas the Laplacian
identity and positive-network interpretation already apply to the
intermediate angle structures. Newton's equation is precisely
\begin{equation}
 2L_{II}(\alpha(u))\,\delta u=\log H(u).
 \label{eq:ipolyNewton}
\end{equation}


% BEGIN INLINE FIGURE: 03_cotangent_hessian.tex
\begin{figure}[htbp]
\centering
\begin{tikzpicture}[font=\small,line cap=round,>=Latex]
  \begin{scope}[shift={(-5.3,0)}]
    \coordinate (fi) at (0,0); \coordinate (fj) at (4,0);
    \coordinate (fk) at (1.5,2.6); \coordinate (fl) at (2.4,-2.4);
    \coordinate (fc1) at (2,.5788461538); \coordinate (fc2) at (2,-.4);
    \draw[mpblue!15] (fc1) circle (2.082081);\draw[mpblue!15] (fc2) circle (2.039608);
    \fill[mpblue!5] (fi)--(fj)--(fk)--cycle;
    \fill[mpteal!5] (fi)--(fl)--(fj)--cycle;
    \draw[mpgray!85] (fi)--(fk)--(fj) (fi)--(fl)--(fj);
    \draw[mpblue,line width=1.3pt] (fi)--(fj);
    \draw[mpochre,line width=1.7pt] (fc1)--(fc2);
    \draw[mpgray!60] (2.12,0)--(2.12,.12)--(2,.12);
    \pic[draw=mpgray,"$\alpha_e$",angle radius=7mm,angle eccentricity=1.5]
      {angle=fi--fk--fj};
    \pic[draw=mpgray,"$\beta_e$",angle radius=7mm,angle eccentricity=1.5]
      {angle=fj--fl--fi};
    \foreach \p in {fi,fj,fk,fl}\fill[mpblue] (\p) circle (1.7pt);
    \fill[mpochre] (fc1) circle (1.7pt) node[right=3pt] {$c_+$};
    \fill[mpochre] (fc2) circle (1.7pt) node[right=3pt] {$c_-$};
    \node[left] at (fi) {$i$};\node[right] at (fj) {$j$};
    \node[above] at (fk) {$k$};\node[below] at (fl) {$l$};
    \node[below,text=mpblue] at (1,0) {$e$};
    \node[left,text=mpochre] at (2,.32) {$e^*$};
  \end{scope}
  \node[align=left,text width=6.3cm,anchor=west] at (.35,.55) {%
    \textbf{One local weight, several roles}\\[7pt]
    $\displaystyle w_e=\frac12(\cot\alpha_e+\cot\beta_e)$\\[6pt]
    $\displaystyle\hphantom{w_e}=\frac{|e^*|_\pm}{|e|}$\\[9pt]
    $\displaystyle\int |\nabla u|^2=\sum_{e=ij}w_e(u_i-u_j)^2$\\[9pt]
    $\displaystyle -D_u^2\mathcal V=2L_{II}$\\[8pt]
    \textcolor{mpgray}{The last identity uses the specified angle
    coordinates and fixed boundary values.}};
\end{tikzpicture}
\caption[The cotangent Hessian and its finite-element interpretation]{The cotangent weight for two adjacent Euclidean triangles.
The dual segment joins their circumcenters and has signed length
$|e^*|_\pm$; in the locally Delaunay configuration shown, its sign is
positive. The finite-element stiffness matrix is the same cotangent
matrix that appears as minus one half of the volume Hessian in the
coordinates $\alpha_c(u)=\alpha_c^0+u_i-u_j$, with fixed edge-angle
sums and $u=0$ on the boundary. This identity is specific to that
coordinate and constraint system.}
\label{fig:mono-cotangent-hessian}
\end{figure}
% END INLINE FIGURE

\subsection{Boundary response, Schur complements, and a Gaussian factor}

For the finite element network, eliminating the interior voltage gives
\begin{equation}
 v_I=-L_{II}^{-1}L_{IB}v_B,\qquad
 \mathcal N=L_{BB}-L_{BI}L_{II}^{-1}L_{IB}.
 \label{eq:ipolyDtN}
\end{equation}
The matrix $\mathcal N$ maps boundary voltages to boundary currents;
it is the discrete Dirichlet-to-Neumann operator. It has the constant
mode in its kernel. Grounding one boundary vertex, or working modulo
constants, removes that mode.

The same elimination appears in volume response. For a nondegenerate
concave maximum, write
\[
 \mathscr V(x)=V_0-\tfrac12x^TKx+O(|x|^3),\qquad K>0,
\]
and let $Px$ be the selected linear boundary data. Perturbation by a
source $j$ gives
\begin{align}
 \max_x\{\mathscr V(x)+j^TPx\}
    &=V_0+\tfrac12j^T\chi j+O(|j|^3),
 &\chi&=PK^{-1}P^T,\notag\\
 \max_{Px=b}\mathscr V(x)
    &=V_0-\tfrac12b^T\chi^{-1}b+O(|b|^3),
 \label{eq:ipolyBoundaryResponse}
\end{align}
where the second formula assumes $P$ has full row rank. The first is
the response to a source; the second fixes its conjugate boundary
variable. They are related by a Legendre transform. In the planar
problem $K=2L_{II}$; in the Dehn filling problem, the same algebra gives
\eqref{eq:cuspHessian} and \eqref{eq:cuspmatrixHessian}. This identifies
an exact common operation without identifying their different spaces
of boundary data.

More generally, split $x=(y,z)$ and eliminate $y$. The effective
quadratic form is
\begin{equation}
 K_{\rm eff}=K_{zz}-K_{zy}K_{yy}^{-1}K_{yz},\qquad
 \det K=\det K_{yy}\det K_{\rm eff}.
 \label{eq:ipolySchur}
\end{equation}
Completing the square gives the exact real Gaussian identity
\begin{equation}
 \int_{\mathbb R^{\dim y}}
 e^{-\frac1{2\hbar}(y,z)^TK(y,z)}\,dy
 =\frac{(2\pi\hbar)^{\dim y/2}}{\sqrt{\det K_{yy}}}
       e^{-z^TK_{\rm eff}z/(2\hbar)}.
 \label{eq:ipolyGaussian}
\end{equation}
Thus elimination in the Dirichlet problem, reduction of the classical
boundary action, and its leading Gaussian integration use the same
Schur complement. An oscillatory or complex phase requires the specified
contour and determinant branch, as in the later quantum calculations.
The finite-dimensional identity does not supply those analytic data
or identify an entire quantum invariant by itself.

There is also a combinatorial formula for the planar determinant. Let
$\mathfrak F_B$ be the spanning forests of the weighted graph in which
each component contains exactly one boundary vertex. Then
\begin{equation}
 \boxed{\det(-D_u^2\mathcal V)
      =2^{|I|}\sum_{F\in\mathfrak F_B}\prod_{e\in F}w_e.}
 \label{eq:ipolyForestDeterminant}
\end{equation}
To prove it, write $L_{II}=C_I\operatorname{diag}(w_e)C_I^T$, where
$C_I$ is the incidence matrix with boundary rows removed. Cauchy--Binet
expands its determinant over $|I|$-edge subsets. An incidence minor is
$\pm1$ exactly for a forest rooted at $B$, and zero otherwise. This
rooted-forest formula expresses the local Gaussian determinant in
geometric conductances. It is a deduction from the Hessian identity
and finite-dimensional linear algebra.

\subsection{Why the nonlinear iteration needs a barrier}

A fast linear solve does not by itself bound the number of Newton
steps. In particular, $-\mathcal V$ is not uniformly self-concordant:
near a vanishing angle, its entropy-like behavior makes the normalized
third derivative unbounded. The companion study supplies a more useful
estimate involving the barrier
\begin{equation}
 F(\alpha)=-\sum_c\log\sin\alpha_c.
 \label{eq:ipolyBarrier}
\end{equation}
The coefficient $\csc^2\alpha_c$ occurs in both $D^2F$ and
$D^3\mathcal V$, but these are tensors of different order. The precise
relation needed for optimization is the following inequality.

\begin{proposition}[A barrier compatible with ideal volume]
\label{prop:ipolyCompatibility}
For every positive angle structure and tangent vector $h$ whose
coordinates sum to zero in each triangle,
\begin{equation}
 |D^3\mathcal V[h,h,h]|
 \le\sqrt3\,(-D^2\mathcal V[h,h])\sqrt{D^2F[h,h]}.
 \label{eq:ipolyCompatibility}
\end{equation}
Consequently $F-t\mathcal V$ is self-concordant with the standard
constant for every $t\ge0$. The barrier $F$ has parameter at most
the number $m$ of corner angles.
\end{proposition}

\begin{proof}
For one triangle set $s_i=\sin\alpha_i$ and choose the cyclic side
directions $\psi=(0,\alpha_1+\alpha_2,\pi+\alpha_2)$.
They satisfy $\sum_i s_i e^{i\psi_i}=0$. Every vector with
$h_1+h_2+h_3=0$ can therefore be written
$h_i=\lambda s_i\cos\varphi_i$, where $\varphi_i=\psi_i-\tau$.
The cotangent energy identity and a trigonometric expansion give
\[
 A:=-D^2\mathcal V[h,h]=\lambda^2s_1s_2s_3,\quad
 B:=D^2F[h,h]=\lambda^2\sum_i\cos^2\varphi_i,
\]
\[
 D^3\mathcal V[h,h,h]
   =\lambda^3s_1s_2s_3\cos(\varphi_1+\varphi_2+\varphi_3).
\]
For the last identity, use $\cos^3 x=(3\cos x+\cos3x)/4$ and the
vanishing side-vector sum. Expanding the cosine of the sum shows
$|\cos(\varphi_1+\varphi_2+\varphi_3)|
\le\sum_i|\cos\varphi_i|\le\sqrt{3\sum_i\cos^2\varphi_i}$.
This proves the inequality on one triangle, and summation proves it
globally since $\sum A_T\sqrt{B_T}\le(\sum A_T)\sqrt{\sum B_T}$.
Each $-\log\sin\alpha_c$ is self-concordant, and, with
$A=-tD^2\mathcal V[h,h]$ and $B=D^2F[h,h]$,
\[
 |D^3(F-t\mathcal V)[h,h,h]|
 \le2B^{3/2}+\sqrt3 A\sqrt B\le2(A+B)^{3/2}.
\]
Finally Cauchy--Schwarz gives $|DF[h]|\le\sqrt m\sqrt{D^2F[h,h]}$,
the required barrier-parameter bound. Restriction to additional affine
angle constraints preserves these inequalities.
\end{proof}

For the central point $\alpha(t)$ minimizing $F-t\mathcal V$, the
volume gap is at most $m/t$. Indeed, concavity of sine gives
$DF(\alpha)[\beta-\alpha]\le m$ for every feasible $\beta$;
combine this with stationarity and concavity of $\mathcal V$.
Increasing $t$ by a factor $1+c/\sqrt m$, with fixed sufficiently
small $c$, gives a bounded Newton decrement: the Hessian of
$F-t\mathcal V$ dominates $D^2F$ and the barrier gradient has local
norm at most $\sqrt m$. Standard self-concordant Newton analysis
therefore gives, starting from a centered point at $t_0>0$,
\begin{equation}
 O\!\left(\sqrt m\log\max\left\{e,\frac{m}{t_0\varepsilon}\right\}\right)
 \quad\hbox{Newton steps to volume gap at most }\varepsilon.
 \label{eq:ipolyStepBound}
\end{equation}
See~\cite[Chapter~11]{BoydVandenberghe} for the underlying path-following
analysis. This is a bound for the barrier method, not for unrestricted
plain Newton iteration.

In the disk coordinates the Hessian of $F-t\mathcal V$ is again a
grounded planar Laplacian, now with interior conductances
\[
 2t w_e+\csc^2\alpha_e+\csc^2\beta_e.
\]
Nearly-linear Laplacian solvers approximate a linear solve in its energy
norm~\cite{SpielmanTengSolvers}. With controlled inexact-Newton
tolerances, their combination with \eqref{eq:ipolyStepBound} gives the
usual $\widetilde O(n^{3/2})$ arithmetic-work scale after centering,
up to accuracy factors, when $m=O(n)$. It does not give a nearly-linear
bound for the whole nonlinear realization problem. Finding a feasible
point, centering it, controlling arithmetic precision, and certifying
the final geometric error are separate costs. The numerical timings in
the companion are experiments on selected inputs, not bounds for these
additional tasks.

\subsection{Flips connect the local matrix to mapping-class geometry}

A diagonal exchange has two different roles. Among triangulations of
the convex hull of a fixed planar vertex set, with fixed nodal values,
Rippa's theorem says that Delaunay
triangulation minimizes the piecewise-linear Dirichlet energy; in matrix
terms $L(\mathcal T_D)\preccurlyeq L(\mathcal T)$~\cite{RippaEnergy}.
This is an ordering of quadratic forms for the same geometric data.
It does not assert monotonicity of hyperbolic volume along an arbitrary
sequence of topological flips.


% BEGIN INLINE FIGURE: 08_flip.tex
\begin{figure}[htbp]
\centering
\begin{tikzpicture}[font=\small,>=Latex,line cap=round,line join=round]
  \begin{scope}[shift={(-5.3,0)}]
    \coordinate (fa) at (0,0);\coordinate (fb) at (3.2,0);
    \coordinate (fc) at (2.7,2);\coordinate (fd) at (.3,1.5);
    \draw[mpblue!15] (1.6,.6625) circle (1.731735);
    \fill[mpblue!5] (fa)--(fb)--(fc)--cycle;
    \fill[mpteal!5] (fa)--(fc)--(fd)--cycle;
    \draw[mpgray] (fa)--(fb)--(fc)--(fd)--cycle;
    \draw[mpochre,line width=1.35pt] (fa)--(fc);
    \foreach \p in {fa,fb,fc,fd}\fill[mpblue] (\p) circle (1.8pt);
    \node[below left] at (fa) {$A$};\node[below right] at (fb) {$B$};
    \node[above right] at (fc) {$C$};\node[above left] at (fd) {$D$};
    \node[font=\bfseries\small] at (1.6,2.85) {The diagonal $AC$};
    \node[text=mpochre,align=center,font=\scriptsize] at (1.6,-1.35)
      {Not locally Delaunay\\$D$ lies inside the circle through $A,B,C$};
  \end{scope}
  \draw[->,mpgray,line width=.85pt] (-1.1,.85)--(.9,.85)
    node[midway,above] {flip};
  \begin{scope}[shift={(1.8,0)}]
    \coordinate (ga) at (0,0);\coordinate (gb) at (3.2,0);
    \coordinate (gc) at (2.7,2);\coordinate (gd) at (.3,1.5);
    \fill[mpblue!5] (ga)--(gb)--(gd)--cycle;
    \fill[mpteal!5] (gb)--(gc)--(gd)--cycle;
    \draw[mpgray] (ga)--(gb)--(gc)--(gd)--cycle;
    \draw[mpteal,line width=1.35pt] (gb)--(gd);
    \foreach \p in {ga,gb,gc,gd}\fill[mpblue] (\p) circle (1.8pt);
    \node[below left] at (ga) {$A$};\node[below right] at (gb) {$B$};
    \node[above right] at (gc) {$C$};\node[above left] at (gd) {$D$};
    \node[font=\bfseries\small] at (1.6,2.85) {The diagonal $BD$};
    \node[text=mpteal,align=center,font=\scriptsize] at (1.6,-1.35)
      {Locally Delaunay\\the planar vertices have not moved};
  \end{scope}
  \node[align=center] at (0,-2.35)
    {For the same nodal values $u$: \quad
     $D_{\mathcal T_{BD}}(u)\le D_{\mathcal T_{AC}}(u)$.};
\end{tikzpicture}
\caption[A diagonal exchange and the Delaunay energy comparison]{A diagonal flip on a fixed Euclidean quadrilateral, shown
with its actual vertex positions. Rippa's energy comparison selects the
Delaunay diagonal for fixed nodal values. The same combinatorial move
also occurs in flip graphs used in surface topology; the displayed
energy comparison concerns the fixed Euclidean realization.}
\label{fig:mono-flip}
\end{figure}
% END INLINE FIGURE

For a connected orientable finite-type marked surface admitting an ideal
triangulation, with at least one marked point on every boundary curve,
the graph of ideal triangulations and diagonal exchanges is instead
a coarse geometric model of its mapping class group. With isotopies
fixing the marked points, the action is proper and cocompact, and an
orbit map is a quasi-isometry~\cite[Section~2.1, Lemma~2.3]{DisarloParlier}.
The pants graph is a different
model, quasi-isometric to Teichm\"uller space with its WP metric, and
Brock relates WP distance to convex-core volume~\cite{BrockWPVolumes}.
His mapping-torus theorem relates WP translation distance and hyperbolic
volume for closed pseudo-Anosov surface bundles~\cite{BrockMappingTori}.
These comparisons have constants depending on the surface topology;
the flip graph and the pants graph should not be identified.

There is a direct volume estimate when a flip sequence gives a layered
ideal triangulation of a hyperbolic mapping torus with $N$ tetrahedra:
straightening gives $\operatorname{Vol}(M)\le Nv_3$, where $v_3$ is the
volume of the regular ideal tetrahedron. This does not require the
given triangulation to have positive geometric shapes; signed
straightened volumes still have absolute value at most $v_3$.
It connects a combinatorial decomposition length to three-dimensional
volume. Comparing it further to WP geometry requires a theorem for the
same surface and boundary conventions, rather than transferring the
closed-surface theorem silently to a punctured surface.

\subsection{The three-manifold extension and the remaining questions}

The Casson--Rivin program uses the same ideal-tetrahedron volume and the
same angle-sum constraints~\cite{FuterGueritaud}. Its local Hessian is
still a cotangent quadratic form. For leading--trailing edge variables
$y_e$, put $p_\alpha=y_e+y_{e'}$ for the pair of opposite tetrahedron
edges carrying $\alpha$, and cyclically for $\beta,\gamma$. Then
\begin{equation}
 -D^2\mathcal V
 =\sum_t\left[
 \cot\alpha_t(dp_\gamma-dp_\beta)^2
 +\cot\beta_t(dp_\alpha-dp_\gamma)^2
 +\cot\gamma_t(dp_\beta-dp_\alpha)^2\right].
 \label{eq:ipolyThreeDimensionalHessian}
\end{equation}
This displays the Euclidean cusp-link triangle inside the second
variation. But a square now involves sums of opposite-edge variables,
and the assembled matrix need not have nonpositive off-diagonal entries
or be diagonally dominant. Additional peripheral directions and removal
of linear dependencies must also be included in a full coordinate
system. The planar positive-Laplacian solver conclusion therefore does
not transfer automatically to arbitrary three-manifold triangulations.
The trianglewise barrier estimate does transfer, because it survives
affine restriction.

An interior critical point has the geometric interpretation supplied by
the Casson--Rivin theorem. A boundary maximizer need not realize the
chosen triangulation by positive ideal tetrahedra. Optimization on a
given angle polytope therefore does not solve the problem of finding
a suitable geometric triangulation or yield a polynomial-time
geometrization algorithm.

The finite element comparison also does not identify its real symmetric
Hessian with the alternating WP form. The exact links established here
are the volume Hessian, its inverse response, its Schur complement,
and its Gaussian determinant. A continuum limit relating appropriately
renormalized polyhedral functionals to a WP potential would require
geometric convergence, a specified normalization, and control of second
variations. This is a natural question, not a theorem asserted here.
Likewise, in three dimensions it remains useful to ask for an effective
bound on the smallest angle at an interior maximum in terms of the
number of tetrahedra and a quantitative feasible margin. Such a bound
would connect variational existence to numerical certification and
complexity. These questions describe how the computational strand can
extend the monograph's variational and semiclassical program.

